\begin{theorem}[G-17: Exact criterion for four-row stretched directions]\label{G-17}
Assume $n>m\geq3$, put $a=n-m$, and suppose
\[
 \dim\mathcal P_{m,4}+3>2n^2+2.
\]
Fix $s\geq1$ and a partition $\nu\vdash ns$ of length at most four,
appending zero parts when necessary. There is a positive integer $j$
such that
\begin{equation}\label{G17:gap}
 m_{j\nu}(\ell^a\operatorname{per}_m)>m_{j\nu}(\det_n),
 \qquad\text{in coordinate degree }sj,
\end{equation}
if and only if
\begin{equation}\label{G17:criterion}
 \nu_4>0\qquad\hbox{and}\qquad\nu_1>as.
\end{equation}
If either condition fails, no positive stretch $j$ gives a gap.
If both hold, an explicit progression of stretches gives gaps at
every sufficiently large term, as specified below. This does not
assert a gap at every sufficiently large integer stretch.

Here are effective bounds in the affirmative case. Put
\[
 \eta=(\nu_1-\nu_4,\nu_2-\nu_4,\nu_3-\nu_4,0),\qquad e=|\eta|,
 \qquad r_\eta=\#\{(i,j):1\leq i<j\leq4,\ \eta_i>\eta_j\},
\]
and define
\[
 \begin{aligned}
 b&=\binom{\binom{n+3}{3}+s-1}{s}\dim S_\eta\mathbb C^4,
 &u&=2n^2+r_\eta-13,\\
 L&=ns+3e,&p&=ns+e,\qquad \beta=6bL^{30}.
 \end{aligned}
\]
The Schur dimension is explicitly
$\prod_{i<j}(\eta_i-\eta_j+j-i)/(j-i)$.
One successful stretch satisfies
\begin{equation}\label{G17:firststretch}
 1\leq j\leq
 j_*:=\beta\bigl[(u+1)((p\beta^2)^u-1)+1\bigr].
\end{equation}
For a progression, let
\[
 T=\binom{b+\beta}{b}-1,\quad E=4T(\beta!),\quad
 Q=(pE)^{u-1},\quad K=u(Q-1)+1.
\]
Then \eqref{G17:gap} holds for every stretch $j=Ek$ with $k\geq K$.
\end{theorem}

\paragraph{Meaning of the criterion.}
For any one of the dimension-certified padded pairs, this determines
exactly which rational directions in the four-row partition simplex
contain a stretched multiplicity obstruction. Its two excluded faces
come from low-row determinant survival and the padding weight bound.
In particular this is a classification of all such directions,
not only a further special choice of the one-longer-row labels in G-16.
The coordinate degrees and stretch factors remain large; no efficient
multiplicity-positivity algorithm is asserted.

\begin{proof}
Necessity follows from earlier ideal containments and includes equality.
If $\nu_4=0$, every $j\nu$ has at most three rows, and G-05 gives
the full ambient multiplicity on the determinant side. If
$\nu_1\leq as$, then in coordinate degree $d=sj$ one has
$(j\nu)_1\leq ad$. The coefficient-weight argument of G-14 applies
and gives the opposite weak multiplicity inequality for every $j$.
Thus it remains to prove sufficiency under \eqref{G17:criterion}.

Use $V=\mathbb C^4$ and the two degree-$n$ cones $Y_D,Y_P$ of G-16.
We first give the auxiliary flag construction, including its coordinate
ring, so that no projective-normality or multiplicity-free assertion
is left implicit.

\paragraph{1. The auxiliary cone and its Cartan coordinate ring.}
Let $v_\eta$ be a nonzero extremal vector of weight $-\eta$ in
$S_\eta V^*$. Concretely, put
$c_i=\eta_i-\eta_{i+1}$ for $i=1,2,3$ and take the product of
the vectors
\[
 (e_0^*\wedge\cdots\wedge e_{i-1}^*)^{\otimes c_i}.
\]
It lies in the single Cartan component $S_\eta V^*$ of the indicated
tensor product of exterior powers. Its projective orbit records the
nested subspaces of dimensions $i$ for which $c_i>0$.
This is the corresponding partial flag variety, a closed incidence
subvariety of a product of Grassmannians. Its dimension is
$r_\eta$: its stabilizer is the block parabolic for the equal-part
blocks of $\eta$, and the missing off-diagonal blocks contribute
one parameter for each pair counted in $r_\eta$.
If $\eta=0$, the projective orbit is a point and $r_\eta=0$.

Let $\mathcal C_\eta$ be the affine cone over this projective orbit.
Its degree-$d$ coordinate ring is
\begin{equation}\label{G17:cartanring}
 \mathbb C[\mathcal C_\eta]_d=S_{d\eta}V.
\end{equation}
Here is a direct representation-theoretic proof. The dual of its
degree-$d$ coordinate space is the span of the group translates of
$v_\eta^{\otimes d}$. This is an extremal vector of weight $-d\eta$.
Its weight space at the extreme of the tensor-product weight order
is one dimensional, and it generates exactly the irreducible Cartan
component $S_{d\eta}V^*$. Taking the dual proves
\eqref{G17:cartanring}. Products give its usual Cartan multiplication.
When $\eta=0$, the cone is a line and each graded piece is
one dimensional, so the same formula holds.

For either cone $Y$, define the enhanced Veronese cone
\begin{equation}\label{G17:enhancement}
 X_Y=\{F^{\otimes s}\otimes v:F\in Y,\ v\in\mathcal C_\eta\}^{\rm cl}
 \subseteq W:=\operatorname{Sym}^s(\operatorname{Sym}^n V)
                      \otimes S_\eta V^*.
\end{equation}
Its projectivization is the Segre product of the $s$th Veronese
embedding of $\mathbb P(Y)$ with the partial flag variety. Hence
\begin{equation}\label{G17:dimension}
 \dim X_Y=\dim Y+r_\eta.
\end{equation}
Pullback of degree-$d$ coordinate functions gives exactly
\begin{equation}\label{G17:ring}
 \mathbb C[X_Y]_d=\mathbb C[Y]_{sd}\otimes S_{d\eta}V.
\end{equation}
The first factor is the Veronese subring of the standard graded
coordinate ring of $Y$, and the second is \eqref{G17:cartanring}.
Both are generated by their degree-one pieces in these gradings;
products of paired coordinates span their diagonal tensor product.
The pullback into the product coordinate ring is injective, proving
\eqref{G17:ring} without assuming that $Y$ is normal.

Write the actual degree-$sd$ coordinate representations as
$\bigoplus_{\lambda\vdash nsd}m_\lambda(Y)S_\lambda V^*$.
The invariant multiplicity in \eqref{G17:ring} is
\[
 \dim\operatorname{Hom}_{\mathrm{SL}_4}
                   (S_\lambda V,S_{d\eta}V).
\]
It is one if $\lambda=d\eta+t(1,1,1,1)$, and zero otherwise.
The size $nsd$ forces
$4t=(ns-e)d=4\nu_4d$, so $\lambda=d\nu$.
Consequently
\begin{equation}\label{G17:extraction}
 \dim\mathbb C[X_Y]_d^{\mathrm{SL}_4}=m_{d\nu}(Y).
\end{equation}
This is a single-label extraction, with original coordinate degree
$sd$, rather than a sum over different labels or degrees.

\paragraph{2. Stability for every strict direction in the criterion.}
Define four positive rational numbers
\[
 \theta_0=\frac{\nu_1/s-a}{m},\qquad
 \theta_i=\frac{\nu_{i+1}}{sm}\quad(1\leq i\leq3).
\]
Positivity is exactly \eqref{G17:criterion}, and their sum is one.
Let $u_0=x_0^n$ and
$u_i=\sqrt{\binom nm}\,x_0^a x_i^m$ for $i>0$.
These are orthonormal monomial vectors for the tensor-induced metric.
Put
\[
 F=\sum_{i=0}^3\sqrt{\theta_i}\,u_i,
 \qquad w=F^{\otimes s}\otimes v_\eta,
\]
normalizing $v_\eta$ to have norm one. The degree-$m$ factor of $F$
is a Fermat sum with four nonzero coefficients. The binary factorization
and cyclic-pencil construction in G-16 therefore put $F$ in both
$Y_D$ and $Y_P$, so $w$ belongs to both enhanced cones.

The diagonal moment of $F$, before removing the scalar part, is
\[
 a e_0+m\sum_i\theta_i e_i=\nu/s.
\]
The moment of a unit pure tensor power is $s$ times the moment of
the original vector; this follows by applying the Lie operator to
each factor and taking the tensor inner product. Thus the first
factor of $w$ contributes $\nu$. The auxiliary extremal vector
contributes its weight $-\eta$. Their sum is
$\nu-\eta=\nu_4(1,1,1,1)$, whose traceless part is zero.
Off-diagonal moments vanish on $F$ because its monomial weights
differ by $m$ times a root, with $m\geq3$, and on $v_\eta$ because
it is a weight vector. They therefore vanish on the tensor product.

The projective stabilizer of this $F$ has zero Lie algebra by the
mixed-monomial calculation in G-16; all four coefficients remain
nonzero, as required there. A Lie operator stabilizing the line of
$F^{\otimes s}$ must stabilize the line of $F$: in the symmetric
algebra on $\operatorname{Sym}^n V$ the equation
$sF^{s-1}(A\cdot F)=cF^s$ implies $A\cdot F=(c/s)F$.
A projective stabilizer of the decomposable tensor $w$ consequently
stabilizes that same line. Thus its projective stabilizer is finite.
Kempf--Ness gives a stable point of each enhanced cone.

The invariant quotients have dimensions
\begin{equation}\label{G17:quotdims}
 \begin{split}
 \dim A_D&=2n^2+r_\eta-13=u,\\
 \dim A_P&=\dim\mathcal P_{m,4}+r_\eta-12>u,
 \end{split}
 \qquad A_D=\mathbb C[X_{Y_D}]^{\mathrm{SL}_4},\quad
 A_P=\mathbb C[X_{Y_P}]^{\mathrm{SL}_4}.
\end{equation}
Indeed stable loci are nonempty and open, so one subtracts the
common orbit dimension $15$ from \eqref{G17:dimension}.
The original dimension gap is preserved for every flag type.

\paragraph{3. The two polynomial degrees and the effective bounds.}
The dimension of $W$ is $b$. Its group-action matrix coefficients
have degree at most $L=ns+3e$ in the entries of $g\in\mathrm{SL}_4$:
the first factor has degree $ns$, while $S_\eta(g^{-\mathsf T})$
has degree at most $3e$ because each entry of $g^{-\mathsf T}$ is
a cubic cofactor. The representation has finite kernel, since its
kernel is contained in the finite projective stabilizer of $w$.
Derksen's Theorem~1.1 and Proposition~1.2 therefore give generators
of $\mathbb C[W]^{\mathrm{SL}_4}$ of degrees at most
$\beta=6bL^{30}$, exactly as in G-16. Exactness of invariants makes
$A_D,A_P$ graded quotients with the same bound on their generators.

For the image-degree bound a smaller, different parameter degree
is valid, and the reason is important. Introduce an independent
matrix $B\in\operatorname{End}(V^*)$, without imposing any relation
between $B$ and the determinant-pencil parameters. The map
\[
 B\longmapsto S_\eta(B)v_\eta
\]
is a homogeneous polynomial map of degree $e$. One can see this
directly in the tensor construction of the Schur functor: it is
the restriction of $B^{\otimes e}$ to the relevant Schur component.
The invertible $B$ form a dense open set and give the full
$\mathrm{GL}(V^*)$ orbit, which is dense in $\mathcal C_\eta$
when $e>0$. Its scalar matrices supply every nonzero scalar
multiple, because the base field is $\mathbb C$.
Thus
\begin{equation}\label{G17:parameterization}
 (Z,B)\longmapsto
  \det(Z_0x_0+\cdots+Z_3x_3)^{\otimes s}
       \otimes S_\eta(B)v_\eta
\end{equation}
has dense image in $X_{Y_D}$ and is homogeneous of degree
$p=ns+e$ in its $4n^2+16$ independent parameters.
No inverse-matrix constraint or adjugate is needed for this
parametrization. If $e=0$, the second factor is constant and can
be omitted; all scalar multiples are already supplied by scaling
the first cone. Formula \eqref{G17:parameterization} then remains
valid with degree $p=ns$.

Apply the cumulative generator-filtration comparison from G-16,
now with quotient dimension $u$, generator bound $\beta$, and
parameter degree $p$. For clarity, its two inequalities become
\[
 \sum_{d=0}^{D}\dim(A_D)_d
       \leq(p\beta)^u\binom{D+u}{u},\qquad
 \sum_{d=0}^{\beta h}\dim(A_P)_d
       \geq\binom{h+u+1}{u+1}.
\]
The first uses the ordinary projective closure in the generator
embedding, of degree at most $(p\beta)^u$, and Chardin's theorem.
The second uses $u+1$ algebraically independent generator images.
The same exact binomial ratio is at least
$(h+u+1)/((u+1)\beta^u)$. Taking
$h=(u+1)((p\beta^2)^u-1)+1$ yields a strict cumulative gap through
$j_*=\beta h$. Since degree zero agrees, some positive stretch
$j\leq j_*$ has a strict invariant-dimension gap. Equation
\eqref{G17:extraction} proves \eqref{G17:gap} at that stretch.

For the individual progression, there are at most
$T=\binom{b+\beta}{b}-1$ homogeneous invariant generators.
The explicit grouping argument of G-11 makes the $E$th Veronese
rings standard graded for $E=4T(\beta!)$.
Their dimensions are $u$ and at least $u+1$, and composing their
degree-one invariant coordinates with \eqref{G17:parameterization}
gives parameter degree $pE$. Consequently their degree-$k$ Hilbert
functions satisfy
\[
 \dim(A_D)_{Ek}\leq Q\binom{k+u-1}{u-1},\qquad
 \dim(A_P)_{Ek}\geq\binom{k+u}{u},\qquad Q=(pE)^{u-1}.
\]
The ratio is strictly greater than one once
$k\geq u(Q-1)+1$. Applying \eqref{G17:extraction} gives the
specified successful multiples $j=Ek$.
On this progression the determinant-side multiplicity is positive as
well. The positive-dimensional graded domain $A_D$ has a nonzero
positive-degree generator $g$ of some degree $e_0\leq\beta$.
Since $e_0\mid E$, the elements $g^{Ek/e_0}$ are nonzero for every
$k\geq1$. Equation \eqref{G17:extraction} therefore gives
$m_{Ek\nu}(\det_n)\geq1$. Thus the displayed strict gaps compare
two positive multiplicities and are not occurrence gaps.
Finally the four-variable highest-weight restriction transfers
all the multiplicities back to the full orbit closures.
\end{proof}

\paragraph{Consistency with the earlier boundaries.}
If $\nu$ is rectangular, the strict inequality in
\eqref{G17:criterion} is $n/4>a$, equivalently $n<4m/3$,
recovering G-11/G-12 whenever the dimension hypothesis holds.
For the one-longer-row direction in G-16, it holds for all padding
because the extra first-row component compensates for $a$.
At equality $\nu_1=as$, every stretch remains excluded, rather
than merely failing this stability construction.
G-15 further excludes a finite initial segment when
$j(\nu_1-as)\leq3$; this is consistent with the eventual
progression in the strict interior.

